\documentclass[12pt, a4paper]{article}
\usepackage[utf8]{inputenc}
\usepackage[ngerman]{babel}
\usepackage{amsmath, amssymb}
\usepackage{pgfplots}
\pgfplotsset{compat=1.18}
\usepackage{geometry}
\geometry{margin=2.5cm}
\usepackage{parskip}

\title{\textbf{Aufgabe 10 -- Lösung} \\ Methoden I: Mathematik, 2026S}
\author{}
\date{}

\begin{document}
\maketitle

\section*{Angabe}

Gegeben ist die Funktion $f : \mathbb{R} \to \mathbb{R}$ mit $x \mapsto -(x - 3)^2 + 4$.

\begin{itemize}
    \item Skizzieren Sie den Graph der Funktion.
    \item Bestimmen Sie das Bild der Funktion.
    \item Geben Sie das Monotonieverhalten der Funktion an.
\end{itemize}

% ============================================================
\section*{Was ist das für eine Funktion?}

$f(x) = -(x - 3)^2 + 4$ ist eine \textbf{nach unten geöffnete Parabel}.

Das Minus vor $(x-3)^2$ dreht die Parabel um -- sie hat jetzt einen \textbf{höchsten} Punkt statt einen tiefsten.

\textbf{Scheitelform:} $f(x) = a(x - h)^2 + k$ mit $a = -1$, $h = 3$, $k = 4$.

\begin{center}
\fbox{%
    \parbox{0.7\textwidth}{%
        \centering
        $a = -1 < 0$ $\Rightarrow$ Parabel öffnet \textbf{nach unten}

        \vspace{0.3em}
        \textbf{Scheitel (höchster Punkt) = $(3,\, 4)$}
    }%
}
\end{center}

Wenn man ausmultipliziert:
\[
    f(x) = -(x-3)^2 + 4 = -(x^2 - 6x + 9) + 4 = -x^2 + 6x - 9 + 4 = -x^2 + 6x - 5
\]

\textbf{Nullstellen} (wo schneidet die Parabel die $x$-Achse?):
\begin{align*}
    -(x-3)^2 + 4 &= 0 \\
    (x-3)^2 &= 4 \\
    x - 3 &= \pm 2 \\
    x &= 1 \quad \text{oder} \quad x = 5
\end{align*}

% ============================================================
\section*{Teil 1: Graph der Funktion}

Wertetabelle:

\begin{center}
\begin{tabular}{c|ccccccc}
    $x$ & 0 & 1 & 2 & \textbf{3} & 4 & 5 & 6 \\
    \hline
    $f(x)$ & $-5$ & $0$ & $3$ & \textbf{4} & $3$ & $0$ & $-5$
\end{tabular}
\end{center}

\textbf{Beobachtung:} Die Werte sind \textbf{symmetrisch} um $x = 3$. Die Parabel schneidet die $x$-Achse bei $x = 1$ und $x = 5$.

\begin{center}
\begin{tikzpicture}
\begin{axis}[
    axis lines = middle,
    xlabel = {$x$},
    ylabel = {$y$},
    xmin = -1, xmax = 7.5,
    ymin = -7, ymax = 6,
    width = 12cm,
    height = 9cm,
    grid = major,
    xtick = {-1,0,...,7},
    ytick = {-6,-4,...,6},
    samples = 100,
    legend pos = south east,
]
    \addplot[blue, thick, domain=-1:7.5] {-(x-3)^2 + 4};
    \addlegendentry{$f(x) = -(x-3)^2 + 4$}

    % Scheitelpunkt
    \addplot[only marks, mark=*, mark size=3pt, red] coordinates {(3,4)};
    \node[red, above right] at (axis cs:3.2,4) {\textbf{Scheitel $(3, 4)$}};

    % Nullstellen
    \addplot[only marks, mark=*, mark size=2.5pt, green!60!black] coordinates {(1,0) (5,0)};
    \node[green!60!black, above left] at (axis cs:0.8,0.3) {$x = 1$};
    \node[green!60!black, above right] at (axis cs:5.2,0.3) {$x = 5$};

    % Wertetabelle Punkte
    \addplot[only marks, mark=o, mark size=2pt, black] coordinates {
        (0,-5) (2,3) (4,3) (6,-5)
    };

    % Symmetrieachse
    \addplot[dashed, gray, thin] coordinates {(3,-7) (3,5.5)};
    \node[gray, above] at (axis cs:3,5.5) {\small Symmetrieachse $x = 3$};

\end{axis}
\end{tikzpicture}
\end{center}

% ============================================================
\section*{Teil 2: Bild der Funktion}

Das \textbf{Bild} (Wertebereich) ist die Menge aller $y$-Werte, die die Funktion annehmen kann.

\textbf{Überlegung Schritt für Schritt:}

\begin{enumerate}
    \item $(x - 3)^2 \geq 0$ für alle $x \in \mathbb{R}$ \quad (Quadrate sind nie negativ)
    \item Also: $-(x-3)^2 \leq 0$ \quad (Minus dreht das Vorzeichen um!)
    \item Daher: $f(x) = \underbrace{-(x-3)^2}_{\leq\, 0} + 4 \;\leq\; 4$ \quad für alle $x$
    \item Der Wert $4$ wird tatsächlich erreicht, bei $x = 3$: \quad $f(3) = -0 + 4 = 4$ \checkmark
    \item Nach unten gibt es keine Grenze: für immer größere $|x - 3|$ wird $f(x)$ beliebig klein
\end{enumerate}

\begin{center}
\fbox{%
    \parbox{0.5\textwidth}{%
        \centering
        \textbf{Bild von $f$:} $\quad f(\mathbb{R}) = (-\infty,\, 4]$

        \vspace{0.3em}
        D.h. alle reellen Zahlen $\leq 4$.
    }%
}
\end{center}

\textbf{Vergleich mit der alten Aufgabe:} Bei $+(x-3)^2 + 4$ war das Bild $[4, +\infty)$ (alles $\geq 4$). Durch das Minus wird es $(-\infty, 4]$ (alles $\leq 4$). Das Minus dreht alles um!

% ============================================================
\section*{Teil 3: Monotonieverhalten}

\textbf{Achtung: Umgekehrt wie bei einer nach oben offenen Parabel!}

\begin{itemize}
    \item \textbf{Links vom Scheitel} ($x < 3$): Der Graph steigt $\Rightarrow$ \textbf{monoton steigend}

    \textit{Beispiel:} $f(0) = -5 < f(1) = 0 < f(2) = 3 < f(3) = 4$ \checkmark

    \item \textbf{Rechts vom Scheitel} ($x > 3$): Der Graph fällt $\Rightarrow$ \textbf{monoton fallend}

    \textit{Beispiel:} $f(3) = 4 > f(4) = 3 > f(5) = 0 > f(6) = -5$ \checkmark
\end{itemize}

\begin{center}
\fbox{%
    \parbox{0.7\textwidth}{%
        \centering
        $f$ ist \textbf{monoton steigend} auf $(-\infty,\, 3]$

        \vspace{0.3em}
        $f$ ist \textbf{monoton fallend} auf $[3,\, +\infty)$
    }%
}
\end{center}

\textbf{Merkhilfe:}
\begin{itemize}
    \item $a > 0$ (nach oben offen): erst fallend $\searrow$, dann steigend $\nearrow$ (Tal)
    \item $a < 0$ (nach unten offen): erst steigend $\nearrow$, dann fallend $\searrow$ (Berg)
\end{itemize}

\end{document}
